% year: תשפ"ג
% semester: א
% moed: ג
% source: exams/מבחנים/תשפג/מועד ג תשפג.pdf
% number: 2ב
% categories: derivatives

\begin{question}
חשבו את הנגזרת בנקודה $x=0$ של הפונקציה $y(x)$ הנתונה בצורה סתומה על ידי המשוואה
\[ y+x^3y^3-3xy=x^5 \]
\end{question}

\begin{hint}
מצאו תחילה את $y(0)$ מתוך המשוואה, ואז גזרו את שני האגפים לפי $x$ (כלל השרשרת ונגזרת מכפלה).
\end{hint}

\begin{solution}
\textbf{ערך הפונקציה:} נציב $x=0$ במשוואה: $y+0-0=0$, ולכן $y(0)=0$.

\textbf{גזירה סתומה:} נגזור את שני האגפים לפי $x$ (כאשר $y=y(x)$):
\[ y'+3x^2y^3+3x^3y^2y'-3y-3xy'=5x^4. \]
נציב $x=0,\ y=0$:
\[ y'(0)+0+0-3\cdot 0-0=0 \quad\Longrightarrow\quad \boxed{y'(0)=0}. \]
(באופן כללי $y'=\dfrac{5x^4-3x^2y^3+3y}{1+3x^3y^2-3x}$, והמכנה שווה $1\neq 0$ בנקודה $(0,0)$.)
\end{solution}
